\documentclass[12pt]{article}
\usepackage{amsmath, amssymb}
\usepackage{geometry}
\geometry{a4paper, margin=1in}
\usepackage{enumitem}

\title{\textbf{CSE400 - Lecture 18} \\ \Large Marginal PMF, Correlation, Covariance, Joint Gaussian Random Variables}
\author{}
\date{}

\begin{document}

\maketitle

\section*{1. Lecture Overview}

\textbf{Topics Covered}
\begin{itemize}[noitemsep]
    \item Marginal PMF
    \begin{itemize}[noitemsep]
        \item Definition
        \item Marginal PMF from Joint PMF Table
        \item Solved Example
    \end{itemize}
    \item Correlation
    \item Covariance
    \begin{itemize}[noitemsep]
        \item Definition and Interpretation
        \item Algebraic Form and Properties
        \item Solved Examples
    \end{itemize}
    \item Pearson's Correlation Coefficient
    \item Joint Gaussian Random Variables
    \begin{itemize}[noitemsep]
        \item Definition
        \item Algebraic Form
        \item Interpretation
    \end{itemize}
\end{itemize}

\textbf{Logical Progression}
\begin{itemize}[noitemsep]
    \item Joint PMF $\rightarrow$ Marginal PMF (summation over variables)
    \item $\rightarrow$ Product Moment $E[XY]$
    \item $\rightarrow$ Centered Measure (Covariance)
    \item $\rightarrow$ Normalized Measure (Pearson Coefficient)
    \item $\rightarrow$ Continuous Joint Distribution (Gaussian Model)
\end{itemize}

\hrulefill

\section*{2. Topic-wise Detailed Notes}

\subsection*{2.1 Marginal PMF}

\textbf{(A) Concept Explanation} \\
For discrete random variables $X$ and $Y$, the joint PMF is defined as:
\[ \boxed{p_{XY}(x,y)=P(X=x,Y=y)} \]
\textbf{Key Idea:} The marginal PMF of one variable is obtained by summing the joint PMF over all possible values of the other variable.

\textbf{(B) Key Formulas} \\
\textit{Fundamental}
\[ \boxed{p_{XY}(x,y)=P(X=x,Y=y)} \]
\textit{Marginalization}
\[ \boxed{p_{X}(x)=\sum_{y}p_{XY}(x,y)} \]
\[ \boxed{p_{Y}(y)=\sum_{x}p_{XY}(x,y)} \]

\textbf{(C) Derivation} \\
From the law of total probability:
\[ P(X=x)=\sum_{y}P(X=x,Y=y) \]
Thus:
\[ p_{X}(x)=\sum_{y}p_{XY}(x,y) \]
Similarly:
\[ p_{Y}(y)=\sum_{x}p_{XY}(x,y) \]

\textbf{(D) Properties and Observations} \\
From the table representation:
\begin{itemize}[noitemsep]
    \item Row sums give $p_{X}(x)$
    \item Column sums give $p_{Y}(y)$
\end{itemize}
\[ \boxed{p_{X}(x_{i})=\sum_{j}p_{ij}} \]
\[ \boxed{p_{Y}(y_{j})=\sum_{i}p_{ij}} \]
Normalization:
\[ \boxed{\sum_{i}p_{X}(x_{i})=1} \]
\[ \boxed{\sum_{j}p_{Y}(y_{j})=1} \]

\textbf{(E) Problem-Solving Framework}
\begin{enumerate}[noitemsep]
    \item Identify joint PMF table
    \item Perform row-wise summation $\rightarrow p_{X}(x)$
    \item Perform column-wise summation $\rightarrow p_{Y}(y)$
    \item Verify total probability equals 1
\end{enumerate}

\subsection*{3. Question-wise Analysis: Marginal PMF Example}

\textbf{Given} \\
Joint PMF table:
\begin{center}
\begin{tabular}{|c|c|c|}
\hline
$x \setminus y$ & 0 & 1 \\
\hline
0 & $1/8$ & $1/4$ \\
\hline
1 & $3/8$ & $1/4$ \\
\hline
\end{tabular}
\end{center}

\textbf{Concept Being Tested} \\
Extraction of marginal PMFs from joint PMF

\textbf{Solution}
\[ \boxed{p_{X}(0)=\frac{1}{8}+\frac{1}{4}=\frac{3}{8}} \]
\[ \boxed{p_{X}(1)=\frac{3}{8}+\frac{1}{4}=\frac{5}{8}} \]
\[ \boxed{p_{Y}(0)=\frac{1}{8}+\frac{3}{8}=\frac{1}{2}} \]
\[ \boxed{p_{Y}(1)=\frac{1}{4}+\frac{1}{4}=\frac{1}{2}} \]

\textbf{Final Answer}
\[ p_{X}(x)=\begin{cases} 3/8, & x=0 \\ 5/8, & x=1 \end{cases} \quad \quad p_{Y}(y)=\begin{cases} 1/2, & y=0 \\ 1/2, & y=1 \end{cases} \]

\textbf{Common Errors}
\begin{itemize}[noitemsep]
    \item Missing a term in summation
    \item Confusing row and column aggregation
\end{itemize}

\hrulefill

\subsection*{2.2 Correlation}

\textbf{(A) Concept Explanation} \\
A basic measure of joint behavior is the product moment:
\[ \boxed{R_{XY}=E[XY]} \]
\textbf{Interpretation:}
\begin{itemize}[noitemsep]
    \item Measures average value of product $XY$
    \item Computed about the origin
\end{itemize}

\textbf{(D) Properties and Observations}
\begin{itemize}[noitemsep]
    \item Depends on scale of $X$ and $Y$
    \item Combines: mean, spread, and dependence
\end{itemize}
\[ \boxed{R_{XY}=0 \implies \text{orthogonal about origin}} \]
\textbf{Conclusion:} Not a clean measure of centered dependence.

\textbf{(E) Transition} \\
To eliminate mean effects $\rightarrow$ define covariance.

\hrulefill

\subsection*{2.3 Covariance}

\textbf{(A) Concept Explanation} \\
Covariance is defined as:
\[ \boxed{Cov(X,Y)=E[(X-\mu_{X})(Y-\mu_{Y})]} \]
\textbf{Interpretation:}
\begin{itemize}[noitemsep]
    \item $Cov(X,Y)>0$: deviations have same sign
    \item $Cov(X,Y)<0$: deviations have opposite sign
    \item $Cov(X,Y)=0$: no linear co-variation
\end{itemize}
\textbf{Important observation:} Magnitude depends on units cannot measure strength directly.

\textbf{(B) Key Formulas} \\
\textit{Derived Form}
\[ \boxed{Cov(X,Y)=E[XY]-E[X]E[Y]} \]

\textbf{(C) Derivation} 
\begin{align*}
Cov(X,Y) &= E[(X-\mu_{X})(Y-\mu_{Y})] \\
&= E[XY-X\mu_{Y}-Y\mu_{X}+\mu_{X}\mu_{Y}] \\
&= E[XY]-\mu_{Y}E[X]-\mu_{X}E[Y]+\mu_{X}\mu_{Y}
\end{align*}
Since $\mu_{X}=E[X]$, $\mu_{Y}=E[Y]$:
\[ Cov(X,Y) = E[XY]-\mu_{X}\mu_{Y} \]

\textbf{(D) Properties} \\
\textit{Symmetry}
\[ \boxed{Cov(X,Y)=Cov(Y,X)} \]
\textit{Variance as Special Case}
\[ \boxed{Cov(X,X)=Var(X)} \]
\textit{Linearity}
\[ \boxed{Cov(aX+b,cY+d)=ac~Cov(X,Y)} \]
\textit{Independence}
\[ \boxed{X \perp Y \implies Cov(X,Y)=0} \]
\textbf{Important Statement:} Zero covariance means uncorrelated, but does not imply independence.

\textbf{(E) Problem-Solving Framework}
\begin{enumerate}[noitemsep]
    \item Use shortcut formula: $Cov(X,Y)=E[XY]-E[X]E[Y]$
    \item Compute expectations
    \item Substitute
    \item Simplify
\end{enumerate}

\subsection*{3. Question-wise Analysis: Covariance Example 1}

\textbf{Given} \\
$Y=-2X+3$

\textbf{Concept Being Tested} \\
Covariance under linear transformation

\textbf{Solution}
\[ \boxed{Cov(X,Y)=E[XY]-E[X]E[Y]} \]
\textbf{Step 1:}
\[ XY = X(-2X+3) = -2X^{2}+3X \]
\textbf{Step 2:}
\[ E[Y] = -2E[X]+3 \]
\textbf{Step 3:}
\begin{align*}
Cov(X,Y) &= E[-2X^{2}+3X]-E[X](-2E[X]+3) \\
&= -2E[X^{2}]+3E[X]+2(E[X])^{2}-3E[X] \\
&= -2E[X^{2}]+2(E[X])^{2}
\end{align*}
Using:
\[ \boxed{Var(X)=E[X^{2}]-(E[X])^{2}} \]
\textbf{Final:}
\[ \boxed{Cov(X,Y)=-2~Var(X)} \]

\textbf{Key Insight} \\
Covariance scales linearly with transformation coefficients.

\hrulefill

\subsection*{2.4 Pearson's Correlation Coefficient}

\textbf{(A) Concept Explanation} \\
Pearson's coefficient measures strength and direction of linear relationship:
\[ \boxed{\rho_{XY}=\frac{Cov(X,Y)}{\sqrt{Var(X)Var(Y)}}} \]

\textbf{(D) Properties}
\[ \boxed{-1 \le \rho_{XY} \le 1} \]
\begin{itemize}[noitemsep]
    \item Normalized covariance
    \item Unit-free
    \item Measures both strength and direction
\end{itemize}

\textbf{Interpretation}
\[ \boxed{\rho_{XY}>0 \implies \text{positive linear trend}} \]
\[ \boxed{\rho_{XY}<0 \implies \text{negative linear trend}} \]
\[ \boxed{\rho_{XY}=0 \implies \text{no linear correlation}} \]

\hrulefill

\subsection*{2.5 Example 2}

\textbf{Given}
\[ \boxed{f_{XY}(x,y)=\frac{xy}{96}, \quad 0<x<4, \quad 1<y<5} \]

\textbf{Computations} \\
\textit{Expectations}
\[ \boxed{E[X]=\frac{8}{3}} \]
\[ \boxed{E[Y]=\frac{31}{9}} \]
\[ \boxed{E[XY]=\frac{248}{27}} \]
\textit{Linear Expectation}
\[ \boxed{E[2X+3Y]=\frac{47}{3}} \]
\textit{Variances}
\[ \boxed{Var(X)=\frac{8}{9}} \]
\[ \boxed{Var(Y)=\frac{92}{81}} \]
\textit{Covariance}
\[ \boxed{Cov(X,Y)=0} \]
\textit{Pearson Coefficient}
\[ \boxed{\rho_{XY}=0} \]

\hrulefill

\subsection*{2.6 Joint Gaussian Random Variables}

\textbf{(A) Concept Explanation} \\
A pair $(X, Y)$ is jointly Gaussian if:
\begin{itemize}[noitemsep]
    \item Joint PDF is Gaussian
    \item Marginals are Gaussian
\end{itemize}
\textbf{Applications:} Signal processing, Communication systems, Noisy measurements.

\textbf{(B) Key Formula}
\[ \boxed{f_{XY}(x,y)=\frac{1}{2\pi\sigma_{X}\sigma_{Y}\sqrt{1-\rho^{2}}}\exp\left(\frac{-1}{2(1-\rho^{2})}\left[\left(\frac{x-\mu_{X}}{\sigma_{X}}\right)^{2}+\left(\frac{y-\mu_{Y}}{\sigma_{Y}}\right)^{2}-2\rho\left(\frac{x-\mu_{X}}{\sigma_{X}}\right)\left(\frac{y-\mu_{Y}}{\sigma_{Y}}\right)\right]\right)} \]

\textbf{(D) Parameters}
\begin{itemize}[noitemsep]
    \item $\mu_{X}, \mu_{Y}$: means
    \item $\sigma_{X}, \sigma_{Y}$: standard deviations
    \item $\rho$: correlation coefficient
\end{itemize}

\textbf{(E) Interpretation} \\
If $\rho>0$: Higher values of one variable tend to correspond to higher values of the other.

\hrulefill

\section*{4. Final Revision Section}

\textbf{(A) Complete Formula Sheet}
\[ \boxed{p_{X}(x)=\sum_{y}p_{XY}(x,y)} \]
\[ \boxed{p_{Y}(y)=\sum_{x}p_{XY}(x,y)} \]
\[ \boxed{R_{XY}=E[XY]} \]
\[ \boxed{Cov(X,Y)=E[(X-\mu_{X})(Y-\mu_{Y})]} \]
\[ \boxed{Cov(X,Y)=E[XY]-E[X]E[Y]} \]
\[ \boxed{Cov(X,X)=Var(X)} \]
\[ \boxed{Cov(aX+b,cY+d)=ac~Cov(X,Y)} \]
\[ \boxed{\rho_{XY}=\frac{Cov(X,Y)}{\sqrt{Var(X)Var(Y)}}} \]
\[ \boxed{-1 \le \rho_{XY} \le 1} \]

\textbf{(B) Key Concept Summary}
\begin{itemize}[noitemsep]
    \item Marginal PMF obtained via summation over other variable
    \item Correlation measures raw product moment
    \item Covariance measures centered dependence
    \item Pearson coefficient normalizes covariance
    \item Joint Gaussian models continuous dependence
\end{itemize}

\textbf{(C) Standard Problem Patterns}
\begin{itemize}[noitemsep]
    \item Joint PMF $\rightarrow$ Marginal extraction
    \item Expectation computation (sum/integration)
    \item Covariance using shortcut formula
    \item Linear transformation problems
    \item Gaussian parameter interpretation
\end{itemize}

\end{document}
